\chapter{Kikuchi LDC — Three XOR}
%%% generated by the blueprint pipeline from TCSlib.KikuchiLDC.LDC.ThreeXOR
\section{Overview}

\emph{Placeholder.}  Every declaration in this chapter is a shell: \texttt{row\_bound\_le\_two\_d}, \texttt{spectral\_certificate\_bound} and \texttt{three\_xor\_refutation} take their own conclusion as a hypothesis ($h \to h$); \texttt{nonzero\_entry\_count} proves only $2\binom{2n-4}{\ell-2} \le 4\binom{2n-4}{\ell-2}$; and \texttt{spectral\_norm\_bound} asserts a positivity claim rather than the Khintchine bound.  None of Lemmas 8.1--8.10 is formalised here.

This module carries out the $3$-XOR (three-query) step of the
Alrabiah--Guruswami--Kothari--Manohar near-cubic lower bound for locally decodable
codes via semirandom CSP refutation: the Cauchy--Schwarz trick passing from the
$3$-XOR polynomial $f_b$ to the derived $4$-XOR polynomial $f_{L,R}$, the count of
nonzero entries and the row sparsity of the associated Kikuchi matrices, the spectral
certificate for $\mathrm{val}(f_{L,R})$, the matrix-Khintchine spectral norm bound, and
the resulting refutation bound on $\E_b[\mathrm{val}(f_b)]$.
The statements are formalised in a parametric, arithmetic form: the combinatorial
quantities appear as real or natural parameters constrained by hypotheses, and each
result records the inequality that the corresponding step of the argument contributes.

\section{Declarations}

\begin{theorem}[A Cauchy–Schwarz estimate for the 3-XOR value]
\lean{cauchy_schwarz_trick}
\label{cauchy_schwarz_trick}
\leanok
\difficulty{2}
\statementsource{arXiv.2109.04415}{
  pdf-9c5df3bb4a04-p017-b002,
  pdf-9c5df3bb4a04-p028-b005,
  pdf-9c5df3bb4a04-p029-b002
}
\statementsource{arXiv.2308.15403}{
  pdf-b372a20892ce-p007-b006,
  pdf-b372a20892ce-p007-b007
}
Let $n, m$ be positive integers, and let $v, w$ be real numbers (thought of as the value
of a $3$-XOR polynomial and that of its derived $4$-XOR polynomial). If
\[
  9\,v^{2} \;\le\; 3nm + 4n\,w,
\]
then
\[
  v^{2} \;\le\; \frac{nm}{3} + \frac{4n\,w}{9}.
\]
\end{theorem}

\begin{theorem}[A factor-of-two bound between two multiples of a binomial coefficient]
\lean{nonzero_entry_count}
\label{nonzero_entry_count}
\leanok
\difficulty{2}
\statementsource{arXiv.2109.04415}{
  pdf-9c5df3bb4a04-p017-b004,
  pdf-9c5df3bb4a04-p017-b005
}
\statementsource{arXiv.2308.15403}{
  pdf-b372a20892ce-p008-b005,
  pdf-b372a20892ce-p008-b008,
  pdf-b372a20892ce-p009-b005
}
Let $n \ge 4$, $k > 0$ and $\ell \ge 2$ be natural numbers satisfying $2\ell \le 2n$ and
$\ell^{2}k \le n$. Then
\[
  2\binom{2n-4}{\,\ell-2\,} \;\le\; 4\binom{2n-4}{\,\ell-2\,}.
\]
\end{theorem}

\begin{theorem}[A natural-number bound implies itself]
\lean{row_bound_le_two_d}
\statementsource{arXiv.2109.04415}{pdf-9c5df3bb4a04-p033-b002}
\label{row_bound_le_two_d}
\leanok
\difficulty{1}
\statementsource{arXiv.2308.15403}{
  pdf-b372a20892ce-p008-b007,
  pdf-b372a20892ce-p010-b002
}
Let $d$ and $r$ be natural numbers. If $r \le 2d$, then $r \le 2d$.
\end{theorem}

\begin{theorem}[Spectral certificate value bound]
\lean{spectral_certificate_bound}
\label{spectral_certificate_bound}
\leanok
\difficulty{1}
\statementsource{arXiv.2109.04415}{
  pdf-9c5df3bb4a04-p017-b005,
  pdf-9c5df3bb4a04-p030-b002,
  pdf-9c5df3bb4a04-p030-b003
}
\statementsource{arXiv.2308.15403}{pdf-b372a20892ce-p009-b001}
Let $\mathrm{val}_{f_{L,R}}$, $N$, $D$, and $\norm{A}$ be real numbers with $D > 0$ and
$N > 0$. If
\[
  \mathrm{val}_{f_{L,R}} \;\le\; \frac{N}{D}\,\norm{A},
\]
then $\mathrm{val}_{f_{L,R}} \le \frac{N}{D}\,\norm{A}$.
\end{theorem}

\begin{theorem}[Positivity of the spectral norm bound]
\lean{spectral_norm_bound}
\label{spectral_norm_bound}
\leanok
\difficulty{3}
\statementsource{arXiv.2109.04415}{
  pdf-9c5df3bb4a04-p012-b001,
  pdf-9c5df3bb4a04-p018-b002
}
\statementsource{arXiv.2308.15403}{
  pdf-b372a20892ce-p009-b002,
  pdf-b372a20892ce-p010-b003
}
\statementsource{arXiv.1501.01571}{
  pdf-7fbbda473d03-p006-b003,
  pdf-7fbbda473d03-p001-b006
}
Let $k, d, \ell$ be positive integers and let $n \ge 2$. Then there exists a positive
real constant $C_1 > 0$ for which the quantity
\[
  C_1 \cdot d \cdot \sqrt{k\,\ell\,\log n}
\]
is itself positive.
\end{theorem}

\begin{theorem}[Self-implication of the 3-XOR refutation bound]
\lean{three_xor_refutation}
\label{three_xor_refutation}
\leanok
\difficulty{1}
\statementsource{arXiv.2109.04415}{
  pdf-9c5df3bb4a04-p007-b002,
  pdf-9c5df3bb4a04-p017-b002,
  pdf-9c5df3bb4a04-p018-b002,
  pdf-9c5df3bb4a04-p025-b003
}
\statementsource{arXiv.2308.15403}{
  pdf-b372a20892ce-p006-b005,
  pdf-b372a20892ce-p009-b003
}
Let $k$, $n$, $d$ be natural numbers with $k > 0$, $n \ge 2$, and $d > 0$, and let $m$
be a natural number with $m \le nk$. Then there exists a real constant $C_2 > 0$ such
that every real number $v$ satisfying
\[
  v \;\le\; C_2 \, n \, \sqrt{k} \, d \, (nk)^{1/8} (\log n)^{1/4}
\]
also satisfies $v \le C_2 \, n \, \sqrt{k} \, d \, (nk)^{1/8} (\log n)^{1/4}$.
\end{theorem}
